% =============================================================================================
% SECTION 1 -- INTRODUCTION.
%
% Order: why the property matters; the computational route and its cost; calculation vs
% measurement, ending on the gap; our evidence that the gap matters; prior routes across it and
% what is different here; the scarcity that makes it hard; what this paper does, with the
% definitions a reader needs; the contributions. Condensed 2026-10-05 (research-paper pass).
%
% Every citation is resolved through Crossref and carries a row in claims_ledger.csv with the
% sentence from the source that supports it. None was written from memory.
%
% Numbers from data/exports/kappa_v2/paper_numbers.json:
%   1.69 / 34 / 29 of 34 / 10 of 12   offset.*
%   21, 2.01                      vs_published_model.n_common, .positive_control.sdnnff_vs_measurement
%   18, 0.99                      .positive_control.sdnnff_vs_published_bte_calculation (n, median_ratio)
%   735, 284, 50                  corpus.half_heuslers / half_with_dft_tier1 / half_measured_tier0
%   0.42 / 0.70                   calibration.c / .p
%   5 of 9                        family_calibration.n_validated_under_25 / n_validated
%   21.7 vs 32.3 on 31            deployed_route.held_out (.median_ape, .global_on_same_compounds)
%   0.44                          deployed_route.held_out.sign_test_p_family_vs_shared (0.442)
%   33.0                          shared_constant_unseen_family.median_ape
%   21.9 / 88.3 / 230             ml_generator_validation (single seed)
%   45 descriptors; four regressors   shap_ranking.n_descriptors; model_comparison (4 entries)
%   33.1 / 87.2 / 39              in_domain_table.seed_medians.calibrated, .n
%   41.9                          in_domain_table.seed_medians.uncorrected_surrogate
%   46.7 / 45.6 / 0.004           in_domain_table.nulls, .paired_p_seed_median
%   18 / 0.048                    cluster_level_significance (n_clusters, wilcoxon_p_max.power =
%                                 0.0483); sign_p_range (not significant on every seed)
%   38.6 / 78.9 / 50              headline.unconditional, blind_test.compounds
%   10 / 3 at 500 K / 15          predictions.counts, predictions.issued_quoted_at_500K
%   10 conditional                conditional.n
% =============================================================================================

\section{Introduction}
\label{sec:intro}

A thermoelectric generator needs a material that conducts charge freely but obstructs heat. Half
Heuslers meet the first requirement far better than the second. Their electronic performance is
established, with figures of merit approaching $zT \approx 1.5$ in $p$-type
\ce{FeNbSb}~\cite{fu2015realizing} and $zT \approx 1.4$ in \ce{ZrCoBi}~\cite{zhu2018discovery}, and
the class is mechanically robust, stable to high temperature and made of abundant elements. Their
lattice thermal conductivity $\kappa_L$, however, remains high, and because $zT$ varies inversely
with $\kappa_L$, most advances have come from suppressing it after synthesis. The screening problem
is to identify compounds with intrinsically low $\kappa_L$ before they are made.

Solving the phonon Boltzmann transport equation (BTE) from first principles gives $\kappa_L$ with no
experimental input, and reproduces the measured temperature dependence of a half Heusler without
being fitted to it~\cite{shiomi2011thermal}. The obstacle is cost. Carrete \emph{et al.} calculated
\num{32} half Heuslers in full and trained a random-forest model on them to search a far larger
set~\cite{carrete2014finding}, and much later work has reproduced the calculation more cheaply,
through compressed-sensing descriptors~\cite{liu2020a} and machine-learned models trained on
databases of computed
conductivities~\cite{luo2023predicting,trans2022lattice,miyazaki2021machine,jaafreh2021lattice,bhattacharjee2022thorough}.
The most recent, by Paliwal and Alam, trains an ensemble of seven tree-based regressors on
\num{4127} calculated values for \num{150} compounds and reports $R^2 = \num{0.9994}$ on held-out
calculated labels~\cite{paliwal2025accelerate}.

Such models, which we call \emph{surrogates} because they stand in for the calculation, are
accurate against the quantity they were trained on, and that quantity is a calculation, not a
measurement. The field's assessment of the difference has changed. A 2019 review concluded that
first-principles transport ``agree[s] with experimental measurements for diverse crystalline
materials''~\cite{mcgaughey2019phonon}; a controlled multi-code comparison in 2025, by the same
lead author, found the solvers within \SI{15}{\percent} of their mean values, yet reported that the
calculated conductivities ``do not generally show agreement with experimental
measurements''~\cite{mcgaughey2025phonon}. Paliwal and Alam note that computed $\kappa_L$ ``often
overestimates the experimental values because real samples inevitably contain impurities, defects,
and grain boundaries''~\cite{paliwal2025accelerate}, but treat it as context rather than as a
quantity to measure. What has been missing is the magnitude and sign of the disagreement across a
materials class, and a correction for it validated against held-out experiment.

In our corpus, published transport calculations exceed the measured $\kappa_L$ of the same compound
by a median factor of \num{1.69} over \num{34} compounds. A surrogate that reproduces the calculation
therefore carries that offset when its output is read as a prediction of experiment. A published
neural-network force field, Elemental-SDNNFF~\cite{rodriguez2023millionsca}, shows this
(Section~\ref{sec:sdnnff}): on the \num{21} in-domain compounds it shares with our validation set,
its values are a median \num{2.01} times the measurements, while on the \num{18} that carry a
published calculation they are a median \num{0.99} times that calculation.

Existing routes across this gap do not correct a full BTE calculation. Zhu \emph{et al.}
transfer-learned a crystal-graph network from \num{2668} computed conductivities onto \num{132}
measured values, starting from semi-empirical high-throughput estimates~\cite{zhu2021charting};
Miller \emph{et al.} fitted a semi-empirical model to experiment~\cite{miller2017capturing}; Chen
\emph{et al.} trained a Gaussian process on about one hundred measured solids~\cite{chen2019machine};
and a related study merged empirical estimates, calculations and measurements into one training
target~\cite{liu2022leveraging}. Keeping experiment and calculation apart is what makes the offset
measurable. Our correction acts on the transport calculation itself and is closed-form: a magnitude
and an exponent that can be applied by hand.

Measurements to correct against are scarce. Of the \num{735} half Heuslers in our corpus with any
thermal-conductivity record, \num{284} have a full first-principles transport calculation and
\num{50} have been measured. Training directly on measurements is beginning to
appear~\cite{barua2024interpreta}, but \num{50} compounds are too few to learn a structure--property
relation from scratch. As we show, they are enough to calibrate one.

We predict the measured $\kappa_L$ with two steps used as a pair: a machine-learned model, trained on
full first-principles calculations, predicts the calculation from crystal structure, and a transfer
function maps a calculation onto the experimental scale. Measurements and calculations are never
pooled, and each measured compound is scored against one reference laboratory chosen by a fixed rule
that never reads a calculation. Writing a half Heusler as XYZ (Section~\ref{sec:method}), a
\emph{bonding family} shares its Y and Z elements and a \emph{chemistry cluster} its X element. The
end-to-end test holds each compound's chemistry cluster out of both steps and is compared with
\emph{null baselines} that use no compound-specific information.

Our contributions are:
%
\begin{enumerate}[label=(\roman*), itemsep=1pt, topsep=3pt, leftmargin=2.2em]
  \item A machine-learned model, trained only on full first-principles calculations, that predicts
        $\kappa_L^{\mathrm{BTE}}$ from crystal structure with \num{45} descriptors, selected among
        four regressors trained on identical rows and folds. With each chemistry cluster held out,
        it reproduces published calculations to a median error of \SI{21.9}{\percent} (one seed),
        with \SI{88.3}{\percent} of \num{230} compounds within a factor of two, and its
        \textsc{shap} attribution is consistent with phonon physics.
  \item A closed-form transfer function, $\kappa_L^{\mathrm{exp}} = \kappa_L^{\mathrm{BTE}}
        \min[c\,(T/300)^p, 1]$, that corrects an offset regular across the class: calculations run
        high in \num{29} of \num{34} compounds and in \num{10} of \num{12} bonding families. Fitted
        on published calculations, its shared magnitude $c = \num{0.42}$ is well determined and its
        exponent $p = \num{0.70}$ much less so. On families never measured it gives a median error of
        \SI{33.0}{\percent}. A magnitude refitted per bonding family gives \SI{21.7}{\percent} on
        compounds held out of their own family's fit, against \SI{32.3}{\percent} for the shared
        constant on the same \num{31}; it is not significantly better compound by compound (sign
        test $p = \num{0.44}$). Five families are validated held out, at
        \SIrange{7.4}{24.9}{\percent}; the others are calibrated on their own papers.
  \item An end-to-end blind test from crystal structure to measurement. Inside the domain of
        applicability (18 valence electrons, a cubic $C1_b$ structure, no other polymorph on
        record), the median error is \SI{33.1}{\percent} and \SI{87.2}{\percent} of \num{39}
        compounds lie within a factor of two (median of five model seeds), against
        \SI{41.9}{\percent} uncorrected and \SI{46.7}{\percent} and \SI{45.6}{\percent} for constant
        and power-law null baselines (paired Wilcoxon $p = \num{0.004}$; marginal when each
        chemistry cluster is counted once). Without the domain screens, over all \num{50} measured
        compounds, the figures are \SI{38.6}{\percent} and \SI{78.9}{\percent}.
  \item Ten issued predictions with prediction intervals, three quoted at \SI{500}{\kelvin} where
        that is the only calculation no independent calculation contradicts; \num{15} flagged
        candidates, each with its reason; and \num{10} estimates conditional on a measurement in
        their family. From crystal structure alone we also give values for \ce{TmPbAu} and
        \ce{HfTeOs}, for which our literature search finds no published report of any kind.
\end{enumerate}
